% Part: propositional-logic
% Chapter: syntax-and-semantics
% Section: formulas

\documentclass[../../../include/open-logic-section]{subfiles}

\begin{document}

\olfileid{pl}{syn}{fml}

\olsection{Propositionele \usetoken{P}{formula}}

!!^{formula}s van de propositielogica worden uit
\emph{!!{propositional
    variable}s}\iftag{prvFalse}{\iftag{prvTrue}{,}{ en} de
  propositionele constante~$\lfalse$}{}\iftag{prvTrue}{\iftag{prvFalse}{
    en}{} de propositionele constante~$\ltrue$}{} opgebouwd met
\emph{logische connectieven}.

\begin{enumerate}
\item !!^a{denumerable} verzameling~$\PVar$ van !!{propositional variable}s
  $\Obj p_0$, $\Obj p_1$, \dots
\tagitem{prvFalse}{De propositionele constante voor !!{falsity}~$\lfalse$.}{}
\tagitem{prvTrue}{De propositionele constante voor !!{truth}~$\ltrue$.}{}
\item De logische connectieven:
  \startycommalist
  \iftag{prvNot}{\ycomma $\lnot$ (negatie)}{}%
  \iftag{prvAnd}{\ycomma $\land$ (conjunctie)}{}%
  \iftag{prvOr}{\ycomma $\lor$ (disjunctie)}{}%
  \iftag{prvIf}{\ycomma $\lif$ (!!{conditional})}{}%
  \iftag{prvIff}{\ycomma $\liff$ (!!{biconditional})}{}%
\item Leestekens: (, ) en de komma.
\end{enumerate}

We duiden deze taal van de propositielogica aan met $\Lang L_0$.

\iftag{defNot,defOr,defAnd,defIf,defIff,defTrue,defFalse,defEx,defAll}{%
Naast de hierboven ingevoerde primitieve connectieven gebruiken we ook
de volgende \emph{gedefinieerde} symbolen:
\startycommalist
  \iftag{defNot}{\ycomma $\lnot$ (negatie)}{}%
  \iftag{defAnd}{\ycomma $\land$ (conjunctie)}{}%
  \iftag{defOr}{\ycomma $\lor$ (disjunctie)}{}%
  \iftag{defIf}{\ycomma $\lif$ (!!{conditional})}{}%
  \iftag{defIff}{\ycomma $\liff$ (!!{biconditional})}{}%
  \iftag{defFalse}{\ycomma $\lfalse$ (!!{falsity})}{}%
  \iftag{defTrue}{\ycomma $\ltrue$ (!!{truth})}}{}.

\begin{tagblock}{defNot,defOr,defAnd,defIf,defIff,defTrue,defFalse,defEx,defAll}
\begin{explain}
Een gedefinieerd symbool behoort formeel niet tot de taal, maar wordt
als informele afkorting ingevoerd. Daarmee kunnen we formules afkorten
die bij uitsluitend gebruik van primitieve symbolen vrij lang zouden
worden. Dat is een duidelijk voordeel. Een nog groter voordeel is dat
bewijzen korter worden. Elk primitief symbool moet in bewijzen immers
afzonderlijk worden behandeld. Hoe meer primitieve symbolen er zijn,
des te langer onze bewijzen dus worden.
\end{explain}
\end{tagblock}

% Alternate symbols

\begin{intro}
Wellicht kent u andere termen en symbolen dan de hierboven gebruikte.
In logicaboeken (en door docenten) worden $\sim$, $\neg$ en~!\ vaak
gebruikt voor ``negatie'', en $\wedge$, $\cdot$ en $\&$ voor
``conjunctie''. Veelgebruikte symbolen voor het ``conditioneel'' of de
``implicatie'' zijn
$\rightarrow$, $\Rightarrow$ en $\supset$.
\iftag{prvIff,defIff}{Symbolen voor ``biconditioneel'',
  ``bi-implicatie'' of ``(materiële) equivalentie'' zijn
  $\leftrightarrow$, $\Leftrightarrow$ en
  $\equiv$.}{}
\iftag{prvFalse,defFalse}{Het symbool $\lfalse$ wordt onder meer
  ``onwaarheid'', ``falsum'', ``absurdum'' en ``bottom'' genoemd.}{}
\iftag{prvTrue,defTrue}{Het symbool $\ltrue$ wordt onder meer
  ``waarheid'', ``verum'' en ``top'' genoemd.}{}
\end{intro}

\begin{defn}[Formule]
\ollabel{defn:formulas}
De verzameling~$\Frm[L_0]$ van \emph{!!{formula}s} van de
propositielogica wordt als volgt inductief gedefinieerd:
\begin{enumerate}
\tagitem{prvFalse}{$\lfalse$ is een atomaire !!{formula}.}{}

\tagitem{prvTrue}{$\ltrue$ is een atomaire !!{formula}.}{}

\item Elke !!{propositional variable}~$\Obj p_i$ is een atomaire
  !!{formula}.

\tagitem{prvNot}{Als $!A$ !!a{formula} is, dan is $\lnot !A$
  !!a{formula}.}{}

\tagitem{prvAnd}{Als $!A$ en $!B$ !!{formula}s zijn, dan is $(!A \land
  !B)$ !!a{formula}.}{}

\tagitem{prvOr}{Als $!A$ en $!B$ !!{formula}s zijn, dan is $(!A \lor !B)$
  !!a{formula}.}{}

\tagitem{prvIf}{Als $!A$ en $!B$ !!{formula}s zijn, dan is $(!A \lif !B)$
  !!a{formula}.}{}

\tagitem{prvIff}{Als $!A$ en $!B$ !!{formula}s zijn, dan is $(!A \liff !B)$
  !!a{formula}.}{}

\tagitem{limitClause}{Niets anders is !!a{formula}.}{}
\end{enumerate}
\end{defn}

\begin{explain}
De definitie van !!{formula}s is een \emph{inductieve definitie}. In
wezen construeren we de verzameling !!{formula}s in oneindig veel
stadia. In het beginstadium verklaren we alle atomaire formules tot
formules; dit komt overeen met de eerste gevallen van de definitie,
namelijk die voor
\iftag{prvTrue}{$\ltrue$, }{}%
\iftag{prvFalse}{$\lfalse$, }{}%
$\Obj p_i$. Met ``atomaire !!{formula}'' wordt dus elke !!{formula} van
een van deze vormen bedoeld.

De overige gevallen van de definitie geven regels om uit reeds
geconstrueerde !!{formula}s nieuwe !!{formula}s te construeren. In het
tweede stadium kunnen we daarmee !!{formula}s uit atomaire
!!{formula}s construeren. In het derde stadium construeren we nieuwe
formules uit de atomaire formules en uit de formules die in het tweede
stadium zijn verkregen, enzovoort. Een !!{formula} is precies datgene
wat uiteindelijk in zo'n stadium wordt geconstrueerd.
\end{explain}

Wanneer we een formule $(!B \ast !C)$ schrijven die uit $!B$, $!C$ met
een tweeplaatsig connectief~$\ast$ is geconstrueerd, laten we het
buitenste paar haakjes vaak weg en schrijven we eenvoudig~$!B \ast !C$.

\begin{tagblock}{defNot,defOr,defAnd,defIf,defIff,defTrue,defFalse,defEx,defAll}
\begin{defn}
Formules die met de gedefinieerde operatoren zijn geconstrueerd, worden
als volgt opgevat:

\begin{tagenumerate}{defTrue,defFalse,defNot,defOr,defAnd,defIf,defIff,defEx,defAll}

\tagitem{defTrue}{$\ltrue$ is een afkorting voor
  \iftag{prvFalse}{$\lnot\lfalse$}{$(!A \lor \lnot !A)$ voor een vaste
    atomaire !!{formula}~$!A$}.}{}

\tagitem{defFalse}{$\lfalse$ is een afkorting voor
  \iftag{prvTrue}{$\lnot\ltrue$}{$(!A \land \lnot !A)$ voor een vaste
    atomaire !!{formula}~$!A$}.}{}

\tagitem{defNot}{$\lnot !A$ is een afkorting voor $!A \lif \lfalse$.}{}

\tagitem{defOr}{$!A \lor !B$ is een afkorting voor
  \iftag{prvAnd}{$\lnot(\lnot !A \land \lnot !B)$}{$\lnot !A \lif
    !B$}.}{}

\tagitem{defAnd}{$!A \land !B$ is een afkorting voor
  \iftag{prvOr}{$\lnot(\lnot !A \lor \lnot !B)$}{$\lnot (!A \lif
    \lnot !B)$}.}{}

\tagitem{defIf}{$!A \lif !B$ is een afkorting voor
  \iftag{prvOr}{$\lnot !A \lor !B$}{$\lnot (!A \land \lnot !B)$}.}{}

\tagitem{defIff}{$!A \liff !B$ is een afkorting voor $(!A \lif !B) \land (!B
  \lif !A)$.}{}
\end{tagenumerate}
\end{defn}
\end{tagblock}

\begin{defn}[Syntactische identiteit]
Het symbool $\ident$ drukt syntactische identiteit tussen tekenreeksen
uit. Dat wil zeggen: $!A \ident !B$ geldt dan en slechts dan als $!A$ en
$!B$ even lange tekenreeksen zijn en op elke positie hetzelfde symbool
bevatten.
\end{defn}

Aan weerszijden van het symbool $\ident$ mogen tekenreeksen staan die
door aaneenschakeling zijn verkregen. Zo betekent $!A \ident (!B \lor
!C)$ dat de tekenreeks~$!A$ dezelfde is als de tekenreeks die ontstaat
door, in deze volgorde, een openend haakje, de tekenreeks $!B$, het
symbool $\lor$, de tekenreeks~$!C$ en een sluitend haakje aaneen te
schakelen. In dat geval weten we dat het eerste symbool van $!A$ een
openend haakje is, dat $!A$ de tekenreeks $!B$ als deeltekenreeks bevat
(beginnend bij het tweede symbool), dat die deeltekenreeks door $\lor$
wordt gevolgd, enzovoort.
\end{document}
